% GENERATED FILE. Edit canonical lessons, then run npm run generate:books.
% canonical-source-hash: fnv1a64:0d722a8964c7ef8a
% canonical-lessons: RU-C249-smetana, RU-C249-tvorog, RU-C249-svinina, RU-C249-baranina, RU-C249-vetchina, RU-R249-first-pass-words-from-chapters-241-244, RU-R249-second-pass-words-from-chapters-245-249

\chapter{Sour cream, Cottage cheese, Pork, Lamb, Ham}
\label{ch:ru-sour-cream-cottage-cheese-pork-lamb-ham}
\hlchaptermodality{249}

\begin{chapteropening}
\textbf{By the end of this chapter:} \emph{I can say sour cream, cottage cheese, pork, lamb and ham.}

Complete the last lesson of chapter 249: I can say sour cream, cottage cheese, pork, lamb and ham.
\end{chapteropening}

\section[\texorpdfstring{\ru{сметана} (smetána)}{сметана (smetána)}]{\ru{сметана} (smetána) — sour cream}

\label{lesson:RU-C249-smetana}



\begin{quote}
Before the new one: say the Russian for jeans, then the Russian for shorts.
\end{quote}

\subsection*{You'll want to know: \ru{сметана}}
\textbf{\ru{сметана}} (\emph{smetána}) — \textquotedblleft{}sour cream\textquotedblright{}.

\begin{grammarlens}[title={What you've built}]
One more word for this chapter.
\end{grammarlens}

\subsection*{Your turn}
\begin{itemize}
\raggedright
  \item \emph{Say it:} \emph{smetána}
  \item \emph{Say it:} \emph{smetána}, once more
  \item \emph{Say it:} say \emph{smetána}
  \item \emph{Recall:} say the Russian for jeans, then the Russian for shorts, then say \emph{smetána} again
\end{itemize}

\begin{tcolorbox}[breakable,colback=teal!4,colframe=teal!35!black,title={Before you move on}]
What does \ru{сметана} mean? (\textquotedblleft{}Sour cream\textquotedblright{}.) Say it once more.
\end{tcolorbox}
\section[\texorpdfstring{\ru{творог} (tvoróg)}{творог (tvoróg)}]{\ru{творог} (tvoróg) — cottage cheese}

\label{lesson:RU-C249-tvorog}



\begin{quote}
Before the new one: say the Russian for shorts, then the Russian for sour cream.
\end{quote}

\subsection*{You'll want to know: \ru{творог}}
\textbf{\ru{творог}} (\emph{tvoróg}) — \textquotedblleft{}cottage cheese\textquotedblright{}.

\begin{grammarlens}[title={What you've built}]
One more word for this chapter.
\end{grammarlens}

\subsection*{Your turn}
\begin{itemize}
\raggedright
  \item \emph{Say it:} \emph{tvoróg}
  \item \emph{Say it:} \emph{tvoróg}, once more
  \item \emph{Say it:} say \emph{tvoróg}
  \item \emph{Recall:} say the Russian for shorts, then the Russian for sour cream, then say \emph{tvoróg} again
\end{itemize}

\begin{tcolorbox}[breakable,colback=teal!4,colframe=teal!35!black,title={Before you move on}]
What does \ru{творог} mean? (\textquotedblleft{}Cottage cheese\textquotedblright{}.) Say it once more.
\end{tcolorbox}
\section[\texorpdfstring{\ru{свинина} (svinína)}{свинина (svinína)}]{\ru{свинина} (svinína) — pork}

\label{lesson:RU-C249-svinina}



\begin{quote}
Before the new one: say the Russian for sour cream, then the Russian for cottage cheese.
\end{quote}

\subsection*{You'll want to know: \ru{свинина}}
\textbf{\ru{свинина}} (\emph{svinína}) — \textquotedblleft{}pork\textquotedblright{}.

\begin{grammarlens}[title={What you've built}]
One more word for this chapter.
\end{grammarlens}

\subsection*{Your turn}
\begin{itemize}
\raggedright
  \item \emph{Say it:} \emph{svinína}
  \item \emph{Say it:} \emph{svinína}, once more
  \item \emph{Say it:} say \emph{svinína}
  \item \emph{Recall:} say the Russian for sour cream, then the Russian for cottage cheese, then say \emph{svinína} again
\end{itemize}

\begin{tcolorbox}[breakable,colback=teal!4,colframe=teal!35!black,title={Before you move on}]
What does \ru{свинина} mean? (\textquotedblleft{}Pork\textquotedblright{}.) Say it once more.
\end{tcolorbox}
\section[\texorpdfstring{\ru{баранина} (baránina)}{баранина (baránina)}]{\ru{баранина} (baránina) — lamb (meat)}

\label{lesson:RU-C249-baranina}



\begin{quote}
Before the new one: say the Russian for cottage cheese, then the Russian for pork.
\end{quote}

\subsection*{You'll want to know: \ru{баранина}}
\textbf{\ru{баранина}} (\emph{baránina}) — \textquotedblleft{}lamb (meat)\textquotedblright{}.

\begin{grammarlens}[title={What you've built}]
One more word for this chapter.
\end{grammarlens}

\subsection*{Your turn}
\begin{itemize}
\raggedright
  \item \emph{Say it:} \emph{baránina}
  \item \emph{Say it:} \emph{baránina}, once more
  \item \emph{Say it:} say \emph{baránina}
  \item \emph{Recall:} say the Russian for cottage cheese, then the Russian for pork, then say \emph{baránina} again
\end{itemize}

\begin{tcolorbox}[breakable,colback=teal!4,colframe=teal!35!black,title={Before you move on}]
What does \ru{баранина} mean? (\textquotedblleft{}Lamb (meat)\textquotedblright{}.) Say it once more.
\end{tcolorbox}
\section[\texorpdfstring{\ru{ветчина} (vetchiná)}{ветчина (vetchiná)}]{\ru{ветчина} (vetchiná) — ham}

\label{lesson:RU-C249-vetchina}



\begin{quote}
Before the new one: say the Russian for pork, then the Russian for lamb.
\end{quote}

\subsection*{You'll want to know: \ru{ветчина}}
\textbf{\ru{ветчина}} (\emph{vetchiná}) — \textquotedblleft{}ham\textquotedblright{}.

\begin{grammarlens}[title={What you've built}]
One more word for this chapter.
\end{grammarlens}

\subsection*{Your turn}
\begin{itemize}
\raggedright
  \item \emph{Say it:} \emph{vetchiná}
  \item \emph{Say it:} \emph{vetchiná}, once more
  \item \emph{Say it:} say \emph{vetchiná}
  \item \emph{Recall:} say the Russian for pork, then the Russian for lamb, then say \emph{vetchiná} again
\end{itemize}

\begin{tcolorbox}[breakable,colback=teal!4,colframe=teal!35!black,title={Before you move on}]
What does \ru{ветчина} mean? (\textquotedblleft{}Ham\textquotedblright{}.) Say it once more.
\end{tcolorbox}
\section[(dialogue)]{Words from chapters 241-244 — a first pass}

\label{lesson:RU-R249-first-pass-words-from-chapters-241-244}



\begin{quote}
Say the words for lamb and ham.
\end{quote}

\subsection*{Your turn}
Say these aloud:

\begin{itemize}
\raggedright
  \item a basement, a balcony, a yard, opposite, between
  \item a library, a university, a café, a zoo, a circus
  \item young, elderly, tall, low, ugly
  \item ancient, loud, bright, dark, light
\end{itemize}

\begin{tcolorbox}[breakable,colback=teal!4,colframe=teal!35!black,title={Before you move on}]
A basement? (\textbf{\ru{подвал}}.) A balcony? (\textbf{\ru{балкон}}.) A yard? (\textbf{\ru{двор}}.) Opposite? (\textbf{\ru{напротив}}.) Between? (\textbf{\ru{между}}.) A library? (\textbf{\ru{библиотека}}.) A university? (\textbf{\ru{университет}}.) A café? (\textbf{\ru{кафе}}.) A zoo? (\textbf{\ru{зоопарк}}.) A circus? (\textbf{\ru{цирк}}.) Young? (\textbf{\ru{молодой}}.) Elderly? (\textbf{\ru{пожилой}}.) Tall? (\textbf{\ru{высокий}}.) Low? (\textbf{\ru{низкий}}.) Ugly? (\textbf{\ru{некрасивый}}.) Ancient? (\textbf{\ru{древний}}.) Loud? (\textbf{\ru{громкий}}.) Bright? (\textbf{\ru{яркий}}.) Dark? (\textbf{\ru{тёмный}}.) Light? (\textbf{\ru{светлый}}.)
\end{tcolorbox}
\section[(dialogue)]{Words from chapters 245-249 — a second pass}

\label{lesson:RU-R249-second-pass-words-from-chapters-245-249}



\begin{quote}
Say the words for strange and funny.
\end{quote}

\subsection*{Your turn}
Say these aloud:

\begin{itemize}
\raggedright
  \item strange, funny, wonderful, terrible, excellent
  \item a microwave, a kettle, a bowl, a saucer, a mug
  \item a radio, headphones, a passport, a visa, a document
  \item clothes, a raincoat, a suit, jeans, shorts
  \item sour cream, cottage cheese, pork, lamb, ham
\end{itemize}

\begin{tcolorbox}[breakable,colback=teal!4,colframe=teal!35!black,title={Before you move on}]
Strange? (\textbf{\ru{странный}}.) Funny? (\textbf{\ru{смешной}}.) Wonderful? (\textbf{\ru{прекрасный}}.) Terrible? (\textbf{\ru{ужасный}}.) Excellent? (\textbf{\ru{отличный}}.) A microwave? (\textbf{\ru{микроволновка}}.) A kettle? (\textbf{\ru{чайник}}.) A bowl? (\textbf{\ru{миска}}.) A saucer? (\textbf{\ru{блюдце}}.) A mug? (\textbf{\ru{кружка}}.) A radio? (\textbf{\ru{радио}}.) Headphones? (\textbf{\ru{наушники}}.) A passport? (\textbf{\ru{паспорт}}.) A visa? (\textbf{\ru{виза}}.) A document? (\textbf{\ru{документ}}.) Clothes? (\textbf{\ru{одежда}}.) A raincoat? (\textbf{\ru{плащ}}.) A suit? (\textbf{\ru{костюм}}.) Jeans? (\textbf{\ru{джинсы}}.) Shorts? (\textbf{\ru{шорты}}.) Sour cream? (\textbf{\ru{сметана}}.) Cottage cheese? (\textbf{\ru{творог}}.) Pork? (\textbf{\ru{свинина}}.) Lamb? (\textbf{\ru{баранина}}.) Ham? (\textbf{\ru{ветчина}}.)
\end{tcolorbox}
